\section{Algorithm and Implementation}\label{sec:algorithm}

\subsection{MapReduce formulation}
A MapReduce job applies a user \emph{map} function to every input record to produce
intermediate key--value pairs, groups the pairs by key (the \emph{shuffle}), and
applies a \emph{reduce} function to each group~\cite{dean2008mapreduce}. We chose a
single representation for all three jobs: the value is a numeric vector and the
reduce function is the element-wise sum. \Cref{tab:jobs} gives the keys and
vectors; averages and rankings are computed from the summed vectors afterwards, on
data that is at most a few thousand lines.

\begin{table}[t]
\centering
\caption{Keys and value vectors emitted by the mapper}\label{tab:jobs}
\begin{tabular*}{\columnwidth}{@{\extracolsep{\fill}}lll@{}}
\toprule
Job & Key & Value vector \\
\midrule
\texttt{demand} & pickup zone $|$ hour & $(1)$ \\
\texttt{routes} & pickup zone $|$ drop-off zone & $(1)$ \\
\texttt{fare}   & distance band & $(1,\ f,\ c,\ c\,f,\ c\,t)$ \\
\bottomrule
\multicolumn{3}{@{}p{\columnwidth}@{}}{\footnotesize $f$ = fare, $t$ = tip,
$c=1$ for credit-card trips (TLC records tips only for card payments), else $0$.}
\end{tabular*}
\end{table}

\begin{algorithm}[t]
\caption{Mapper, combiner/reducer and post-processing}\label{alg:mr}
\begin{algorithmic}[1]
\Procedure{Map}{line}
  \State $r \gets$ \Call{Parse}{line} \Comment{skip header and invalid rows}
  \If{$r \neq \bot$} \textbf{emit} $(\mathit{key}(r), \mathit{vec}(r))$ \EndIf
\EndProcedure
\Procedure{Reduce}{$k$, $[v_1,\dots,v_m]$} \Comment{also used as combiner}
  \State \textbf{emit} $(k, \sum_{i} v_i)$ \Comment{element-wise}
\EndProcedure
\Procedure{Post}{$\{(k, s_k)\}$}
  \State \texttt{routes}: sort by $s_k[0]$, keep top $N$
  \State \texttt{fare}: avg fare $= s[1]/s[0]$, tip rate $= s[4]/s[3]$
\EndProcedure
\end{algorithmic}
\end{algorithm}

\textbf{Combiner.} Because addition is associative and commutative, the reducer can
also run on each map task's output before the shuffle. For \texttt{fare} this
collapses millions of map outputs into a few lines per map task, so almost nothing
crosses the network; for \texttt{routes}, with $265^2$ possible keys, the reduction is
smaller. We measure both variants (\cref{sec:results}).

\textbf{Data cleaning.} \textsc{Parse} rejects rows with a malformed field, a zone
outside 1--265, a distance outside $(0,100)$ miles, a fare outside $(0,1000)$\,USD or a
negative tip (\cref{sec:data}).

\subsection{Execution engines}
\begin{itemize}
  \item \textbf{Hadoop Streaming}~\cite{hadoop_streaming,noll_streaming}.
        \texttt{mapper.py} and \texttt{reducer.py} read standard input and write
        tab-separated lines; Hadoop runs one mapper per
        128\,MB HDFS block and sorts map output by key before the reducers.
  \item \textbf{Spark RDD} (on YARN~\cite{spark_yarn}). \texttt{spark\_job.py --api rdd} imports the \emph{same}
        \texttt{parse} and \texttt{map\_record} functions; \texttt{flatMap} is the map
        phase and \texttt{reduceByKey} performs combining, shuffle and reduction.
        Intermediate data stays in memory instead of being written to local disk.
  \item \textbf{Spark DataFrame.} \texttt{--api df} expresses the same filters and
        aggregations in Spark SQL, so they are optimised by Catalyst and executed in
        the JVM without per-record Python calls.
  \item \textbf{Serial baseline.} \texttt{local\_baseline.py} applies the same
        functions in one Python process with an in-memory dictionary---the
        ``single computer'' reference.
\end{itemize}
Because all engines share one parser, any difference in run time comes from the
execution engine, not from the logic. \texttt{scripts/compare\_outputs.py} checks that
all engines give identical results; on the real data they do for every tier,
cluster size and environment (\cref{sec:results}).
